% 负责人：A / hycx233。各模块的细节由对应负责人补充。
\section{任务目标}
本项目按课程实践方案\cite{course}完成已知结构识别与解释、新结构候选发现、
多轨道可视化三个必做任务，选做不同条件下的结构差异定位。
数据条件统一为 WT、$\Delta stpA$ 和 $\Delta hns\Delta stpA$，每个条件使用两个重复。
已知结构类型为 OPCID、CHIN 和 CHID，定义与实验背景参照原研究\cite{paper}。

\section{设计思路}
\subsection{共享输入与两类分析目标}
已实现共用的矩阵读取、坐标转换、窗口提取和数据划分。
当前把 10 bp 数据聚合至 100 bp，截取 24 kb 窗口并生成 $128\times128$ 模型输入；
是否采用附加视角，应由真实窗口效果决定。
共享数据流程已在六个样本上验证，方法和分组统计见第 3 节；模型效果尚未评价。

强度比较保留测序深度归一化后的接触信号；
形态分析使用按基因组距离计算的观察值与期望值之比（O/E），再作统一的数值变换。
不能用每个窗口独立缩放到相同均值的数据判断不同条件的强弱变化。
环状染色体上的跨起点窗口需要首尾拼接，背景距离也要考虑环状结构。

\subsection{模型与评价路线}
已知结构识别采用轻量 CNN，报告准确率、Macro-F1、混淆矩阵与简单基线，
结合显著性图和失败案例说明模型依据。
候选发现采用小型卷积自编码器，以重建误差打分，使用 PCA 和 DBSCAN 整理候选表征。
候选应与已知结构及随机背景对照，并在独立重复的相同位置检查复现性。
不能仅凭分类器低置信度或聚类噪声标签宣称发现了新结构。

正式划分需使同一结构、其重复和空间重叠的输入不跨训练、验证、测试集合。
模型选择使用验证集，测试结果用于最终评价。
任务四比较两个突变条件相对 WT 的变化，并用 WT 两个重复的差异作参照。
这些规则将在对应实验章节用实际参数、对照和结果落实。
